\chapter{总结与展望}

\section{主要工作总结}
本文以融合测角与测距观测的空间目标定位为对象，围绕“模型--算法--实现--评估”展开研究，主要完成了以下工作：
\begin{enumerate}
  \item 建立了统一的加权最小二乘目标函数，处理异构观测与异方差噪声；
  \item 给出了测距与角度观测的梯度、海森矩阵解析表达，并与程序模块对应；
  \item 设计了“几何交会初值 + 非线性迭代求精”的求解流程；
  \item 实现并比较了牛顿法与法方程迭代法，验证两者结果一致性；
  \item 构建了协方差近似、置信椭球与灵敏度分析方法，形成完整精度评估链路。
\end{enumerate}

\section{结论}
根据实验结果与分析，得到以下结论：
\begin{itemize}
  \item 在当前观测配置下，融合模型可稳定给出高精度位置估计；
  \item 牛顿法与法方程法均可收敛到同一最优解，前者局部收敛速度更快；
  \item 法化矩阵逆与海森矩阵逆在量级上相符，均可用于协方差近似；
  \item 垂向不确定性显著高于水平不确定性，主要由几何构型决定；
  \item 几何初值对收敛稳定性与计算效率具有关键作用。
\end{itemize}

\section{不足与改进}
本文仍存在如下不足：
\begin{enumerate}
  \item 实验场景以仿真为主，真实数据验证仍需加强；
  \item 当前“LM 方法”未显式引入阻尼项，极端场景鲁棒性可进一步提升；
  \item 小维显式线性求解适用于课程场景，通用工程中应使用更稳定矩阵分解；
  \item 对异常值仅进行了定性讨论，尚未系统实现鲁棒估计。
\end{enumerate}

\section{未来工作展望}
未来可从以下方向继续展开：
\begin{itemize}
  \item 面向真实观测系统构建端到端数据处理流程；
  \item 引入阻尼牛顿/信赖域方法提高强非线性场景收敛性；
  \item 将静态定位扩展到动态跟踪，结合 EKF/UKF 进行时序融合；
  \item 研究带约束与鲁棒损失的混合优化模型；
  \item 开发可视化评估工具，形成可复用课程与科研平台。
\end{itemize}

\section{结语}
本文工作表明，基于严谨的数学建模和清晰的程序实现，融合测角与测距的空间定位问题可以实现高效、可解释、可扩展的求解。相关方法对后续多传感器融合定位研究具有直接参考价值。
